\section{Method}
\label{sec:method}

\subsection{Framework Overview}

Figure~\ref{fig:framework} shows the pipeline in four stages. In Stage 1, the
\textbf{Delay-Embedded Controlled State} turns judged scores and issued interventions into one
vector per turn. In Stage 2, the \textbf{Linear Operator with an Explicit Action Channel} fits
a one-step map on those vectors with the intervention as its own regressor, so that removing it
changes the fit rather than the model class. In Stage 3, \textbf{Reconstruction-Then-Ridge
Estimation} fits the learned-lifting alternative in two stages, keeping its dynamics half
closed-form; Appendix~\ref{app:estimator} gives it, since it estimates the baseline rather than
the operator this paper is about. Stage 4 reads the fitted operator in one of two ways: \textbf{Budget-Pruned Action
Search} spends a fixed budget against it, and \textbf{Operator Diagnostics as a Pre-Investment
Test} asks whether spending that budget can pay at all. This section answers where in that
pipeline the control-theoretic prior stops carrying its weight.

\subsection{Delay-Embedded Controlled State}

\textbf{Why a window of past readouts, with the action inside the same vector.} A single
judged score is a measurement, not a state: it says where the conversation stands and nothing
about where it is heading, so a map fitted on it can only regress the next score on the current
one. Assumption 1 says a window of past scores carries what the current turn omits, and the
window is how we give the map access to that. The action belongs in the same vector rather than
in a side input for a different reason. Once the intervention sits beside the scores as another
coordinate, removing it leaves a strictly smaller state whose fit can be compared against the
full one on the same rows. A side input, by contrast, could be dropped by rewriting the model
class, and the comparison would then measure the rewrite.
\uline{The delay-embedded controlled state therefore does two jobs: the window gives the
operator the history that one turn withholds, and putting the action inside the same vector
makes withholding it an ablation the same estimator can score.}

\textbf{State construction.} Equation~\eqref{eq:delay-state} of Section~\ref{sec:problem}
gives the state, with $L_y$ readout terms and $L_u$ control terms. A turn contributes a row only
when every lagged term exists and the next turn is observed, so deepening the window shortens
the usable stretch of each trajectory. On the trajectory-tracking tasks the control is the
tracking error against a target constant within a trajectory, which makes it a function of the
state's own leading block; that line cannot separate what the action did from what the state was
already doing, and we report its action-channel cells as undefined rather than zero. One
competing account of the window's value survives this construction: scores may persist because
the item is hard, not because one turn propagates into the next. Section~\ref{sec:experiments}
separates the two by giving the memoryless baseline a per-item intercept.

\subsection{Linear Operator with an Explicit Action Channel}

\textbf{Why a linear map whose action is a separate regressor.} The operator has one job in
this paper, which is to say how much of the next state the present state and the issued
intervention account for. A linear map is the weakest model class that can carry both terms at
once, and a class that cannot beat its own trivial baselines is reporting something about the
data rather than about itself. Assumption 2 asks more than linearity: it asks that the
intervention vary enough on its own that its column is not a rewriting of the state columns.
When that holds, the action has a coefficient of its own, and a fit with the column removed is
the same estimator answering the same question with one fewer input.
\uline{The operator is linear so that a shortfall cannot be blamed on capacity. Its action
enters as a separate regressor so that withholding the action produces a comparison rather than
a different model.}

\textbf{Operator and readout.} Equation~\eqref{eq:beh-dynamics} gives the fitted map. The
lifted state $\eta_t$ is the state passed through a feature dictionary, which is the identity
map in the variant we report, and $A$, $B$, $b$ are the transition, control and intercept
blocks of one ridge least-squares solve. The intervention $v_t$ enters as its own column of
that solve. The readout matrix $C$ comes from a second ridge problem fitted on the same lifted
states, not from the same solve as $A$, $B$ and $b$. One variant replaces the single control
column by the intervention and its product with the current score, which lets the same action
carry a different effect depending on where the trajectory already sits.

\textbf{The withheld-action variant.} Dropping the control column and refitting gives an
operator of the same class on the same rows, and the gap between the two fits in
Equation~\eqref{eq:skill} is what Section~\ref{sec:experiments} reports for every column. The
same construction applies to the learned-lifting estimator of Appendix~\ref{app:estimator}, so
the two model classes are compared on the ablation as well as on the fit.

\subsection{Budget-Pruned Action Search}

\textbf{Why search over sequences, and why the budget prunes instead of penalizing.} An
intervention issued now changes what the next state will be, so the value of acting at this
turn depends on what acting later would have bought, and a rule that scores turns one at a time
cannot see that. Searching over sequences is the smallest construction that can. The budget is
a count of interventions an operator is willing to spend, not a price they are willing to pay,
and the difference decides how it enters: a penalty can always be outweighed by a large enough
predicted gain, which would let the controller overspend whenever it is confident. Pruning the
admissible set cannot be outweighed. The comparison this controller is built to face is the
equal-cost fixed schedule of Section~\ref{sec:problem}, which spends the same count without
reading the conversation.
\uline{The action search therefore ranges over sequences so that the value of acting now
includes what acting later would have bought, and enforces the budget by pruning so that no
predicted gain can buy a step the budget does not allow.}

\textbf{Search objective.} Equation~\eqref{eq:mpc-value} gives the value of a branch, where
$\beta$ is the count of interventions still unspent and $\kappa$ is the penalty charged for
issuing one. The intervention $v_t$ is the action the branch commits to at the current turn.

\textbf{Enumeration and pruning.} The controller enumerates the binary action tree to its
planning horizon, scores each branch by the predicted readout summed along it minus the repeat
penalty, and emits the first action of the best branch. The budget enters at every node: once
the count reaches zero, the admissible set at every deeper node is the null action alone. That
collapse is the entire constraint mechanism, and calling this controller a constrained program,
a Lagrangian relaxation or a projection would describe a method we did not run. With a budget of
one the recursion reduces to spending the allowance now against holding it while the rest of the
horizon stays idle. Appendix~\ref{app:threshold} gives a closed-form threshold for the
interaction variant; it agrees with the search only when both branches leave the second step
idle, so we treat it as a remark and not as the decision rule.

\subsection{Operator Diagnostics as a Pre-Investment Test}

\textbf{Why the identification fit answers the investment question, and answers it early.}
A controller that chooses when to act can only exploit structure the operator has already
found; if withholding the action leaves the fit unchanged, there is no such structure for any
policy to trade on, whatever search runs on top. Assumption 3 states that implication, and its
value is that the fit is available before any closed-loop run: the same data that identifies
the operator answers whether closing the loop can pay. We check one direction. A fit that does
not change when the action is withheld licenses the decision not to invest; a fit that does
change licenses nothing here, because we have no task family where that happened and the closed
loop then paid.
\uline{The identification fit therefore answers the investment question because a controller
can only exploit what the fit can see, and answers it before the closed-loop compute is spent,
in one direction only.}

\textbf{Reachability diagnostics.} Equation~\eqref{eq:ctrb} gives the finite-horizon
controllability matrix $\mathcal{C}_H$ and Gramian $W_H$, and Equation~\eqref{eq:spectral}
gives the spectral radius $\rho(A)$ and the state half-life $t_{1/2}$ it implies. Every
quantity here is a diagnostic of one fit. None is compared across arms or across columns, and
each appears in Section~\ref{sec:experiments} beside a task-level result rather than in place
of one.

\textbf{What the five stages deliver.} The state, the operator and its two estimators exist to
make one contrast well posed, between a fit that sees the action and a fit that does not, while
the action search exists to be measured against an equal-cost schedule on the same task.
\uline{Section~\ref{sec:experiments} asks whether the window pays (Section~\ref{sec:problem}'s
Assumption 1), whether the action channel carries anything (Assumption 2), whether a learned
lifting beats the linear operator, and whether the identification answer predicted the
closed-loop one (Assumption 3).}
